\documentclass{article}

% Language setting
% Replace `english' with e.g. `spanish' to change the document language
\usepackage[english]{babel}

% Set page size and margins
% Replace `letterpaper' with `a4paper' for UK/EU standard size
\usepackage[a4paper,top=2cm,bottom=2cm,left=3cm,right=3cm,marginparwidth=1.75cm]{geometry}

% Useful packages
\usepackage{amsmath}
\usepackage{graphicx}
\usepackage[colorlinks=true, allcolors=blue]{hyperref}

\title{Discrete Math Course}
\author{Wisam Yassar Mahardika}
\date{April 4\textsuperscript{th} 2026}

\begin{document}
\maketitle
\clearpage
\tableofcontents
\clearpage
\section{Implications (converse, inverse, contrapositive) and biconditionals}
\subsection{Implication (conditional statements)}
The \textbf{implication} of propositions $ p $ and $ q $ is denoted $ p \implies q $ and read ``if $ p $ then $ q $`` or ``$ p $ implies $ q $``.\\


When the hypothesis is true, the conclusion must be true for the implication to be true. when the hypothesis is false, the conclusion is true.\\


$ p $ denotes ``It is a holiday`` $ q $ denotes ``The store is closed``. \\


\textit{Truth Table for $ p \implies q $}
\begin{tabular}{|c|c|c|}
\hline
$ p $ & $ q $ & $ p \implies q $ \\
\hline
T & T & T \\ 
T & F & F \\
F & T & T \\
F & F & T \\
\hline 
\end{tabular}
\subsubsection{Converse, Inverse, Contrapositive}
From our implication, $ p \implies q $, we can construct 3 new conditional statements.
\begin{enumerate}
	\item \textbf{Converse: }$ q \implies p $ 
	\item \textbf{Inverse: }$ \neg p \implies \neg q $
	\item \textbf{Contrapositive: }$ \neg q \implies \neg p $
\end{enumerate}


\textit{e.g} It is raining is a sufficient condition for me not going to town.
\begin{itemize}
	\item in an ``if then`` form it looks like this: \textbf{If} it is raining, \textbf{then} I won't go to town. $ p \implies q $
	\item In converse it looks like this: \textbf{If} I don't go to town, \textbf{then} it is raining.
	\item In inverse it looks like this: \textbf{If} it is not raining, \textbf{then} I will go to town.
	\item In contrapositive it looks like this: \textbf{If} I will go to town, \textbf{then} it is not raining.\\
\textbf{Note: }Contrapositive always have the same \textit{truth value} as $ p \implies q $.
\end{itemize}
\subsection{Biconditional}
The \textbf{bicondtional} of proposition $ p $ and $ q $ is denoted $ p \Longleftrightarrow{} q $ and read ``$ p $ if and only if $ q $``.\\


For a biconditional to be true, both propositions must share the same truth value. \\


\textit{Truth Table for $ p \Longleftrightarrow q $}
\begin{tabular}{|c|c|c|}
\hline
$ p $ & $ q $ & $ p \Longleftrightarrow q $ \\
\hline
T & T & T \\ 
T & F & F \\
F & T & F \\
F & F & T \\
\hline 
\end{tabular}


Our biconditional $ p \Longleftrightarrow q $ can also be written as a compound proposition. \\
$ (p \Longleftrightarrow q) \equiv (p \implies q) \land (q \implies p) $ \\


\begin{tabular}{|c|c|c|c|c|c|}
\hline
$ p $ & $ q $ & $ p \implies q $ & $ q \implies p $ & $ (p \implies q) \land (q \implies p) $ & $ p \Longleftrightarrow q $ \\
\hline
T & T & T & T & T & T \\ 
T & F & F & T & F & F \\
F & T & T & F & F & F \\
F & F & T & T & T & T \\
\hline 
\end{tabular}
\end{document}
